\section{Problem Statement}

Artificial Neural Networks (ANNs), especially Multilayer Perceptrons (MLPs), have the ability to learn tasks and solve prediction and classification problems. Using learning algorithms such as Backpropagation, it is possible to adjust the weights of an ANN so that it can address a specific task. MLPs are a common form of ANN, composed of multiple layers of neurons, where the input layer receives the data and the output layer provides the predictions. The intermediate layers, also known as hidden layers, perform non-linear transformations that allow the network to learn complex representations of the data. However, many of the tasks solved in everyday life generate additional information over time. An example of this is the behaviour of a financial time series or weather prediction for a given region. In this context, incremental learning arises as a scarcely explored approach that allows the model to learn new information without the need to retrain the entire model with both old and new data.

In current machine learning models, when a dataset is used to train a specific model, that model is functional only for that dataset and the information it represents. However, when it becomes necessary to incorporate new information into the model, that new data must be collected, added to the existing dataset, and then the entire model must be retrained to integrate the new data.

In this context, if a network trained with an initial dataset \(d_1\) must later be trained with a second dataset \(d_2\), the knowledge acquired from \(d_1\) will be lost. If an incremental learning model is not used, \(d_1\) and \(d_2\) must be combined into a single dataset and the ANN retrained to incorporate the new knowledge (\(d_2\)). On the other hand, if incremental learning is applied, the ANN is first trained with \(d_1\) and then with \(d_2\), maintaining minimal information loss from \(d_1\). If additional information about the problem (\(d_3\)) is obtained in the future, the ANN will be trained with \(d_3\) using the incremental model, minimising the loss of data from \(d_1\) and \(d_2\).

This work is based on the research of \cite{bullinaria2009}, which uses a configuration of duplicated weights in the ANN. In this configuration, the ANN’s weights are duplicated, and different learning rates are assigned to the layers. The first layer, associated with a high learning rate, simulates short-term memory by learning a new task quickly. This research is based on the study of \cite{bullinaria2009}, which proposes a strategy to improve the incremental learning capability of Artificial Neural Networks (ANNs) through weight duplication. In this approach, each neural connection has two versions of its weights: one that simulates short-term memory and another that represents long-term memory.

Short-term memory uses a high learning rate, allowing it to adapt quickly to new information but also making it more prone to forgetting previous knowledge. In contrast, long-term memory is trained with a low learning rate, learning more slowly but retaining previously acquired information more effectively.

During training, both sets of weights are updated in parallel according to their respective learning rates. When making a prediction, the results are obtained from a combination of both sets of weights. This strategy allows the neural network to integrate both recent and accumulated knowledge, emulating the complementary functioning of biological memory systems and forgetting previous information more gradually. The second layer, with a low learning rate, simulates long-term memory, learning more slowly and forgetting less prior information. By integrating both layers, the ANN balances the output to consider both the newly acquired information and the information retained in both weight layers.

The main problem of incremental learning in \cite{bullinaria2009} is that, as new datasets are incorporated, the knowledge acquired from the earlier datasets is progressively lost, which becomes less useful if, in the future, there are 10 or 20 incremental training stages with new data.

Therefore, it is crucial to explore new ANN configurations that improve current methods to reduce the amount of information forgotten as new data arrive. As in previous research, this work will be based on the concepts of short- and long-term memory. However, instead of duplicating the weights and having two learning rates, the idea of creating more copies of the weights will be proposed, each with a distinct learning rate.
